%==============================================================================
%  2. Problem Formulation and DMDE
%==============================================================================

\section{Problem Formulation and DMDE}
\label{sec:problem}

\subsection{Multi-UAV Cooperative Path Planning}
\label{subsec:problem-def}

% 建议内容：决策变量定义 -> 代价函数 -> 约束 -> 与已有模型的关系
% 状态空间：每架无人机的路径表示为若干航路点（经度/纬度/高度）

Consider a fleet of \Np{} unmanned aerial vehicles that must be routed from a
common start region to \Mt{} targets over a terrain represented by a digital
elevation model (DEM). Each path is encoded as a sequence of waypoints whose
coordinates are the decision variables of the problem. The cost of a candidate
solution aggregates three terms:

\begin{equation}
  J(\vx) = w_{1} J_{\mathrm{terrain}}(\vx)
         + w_{2} J_{\mathrm{threat}}(\vx)
         + w_{3} J_{\mathrm{length}}(\vx),
  \label{eq:cost}
\end{equation}

where $J_{\mathrm{terrain}}$ penalises violations of the minimum safe clearance
above the terrain, $J_{\mathrm{threat}}$ quantifies exposure to radar threat
fields, and $J_{\mathrm{length}}$ measures the total flight distance.
% 待补：约束条件（高度上下界、最小离地间隙、单机航程、目标分配唯一性）
\todo{state the constraints: altitude bounds, minimum clearance, flight range
      and one-to-one target assignment.}

\subsection{Differential Evolution Baseline (DMDE)}
\label{subsec:dmde}

DE maintains a population of $\NP$ candidate solutions and iteratively applies
mutation, crossover and selection \citep{storn1997de}. The variant used as our
baseline is referred to as \dmde{}; its mutation uses the current-to-best
strategy with scaling factor \Ffactor{} and crossover probability \CR{}:

\begin{equation}
  \bm{v}_{i} = \bm{x}_{i} + \Ffactor\,(\bm{x}_{\mathrm{best}} - \bm{x}_{i})
             + \Ffactor\,(\bm{x}_{r_{1}} - \bm{x}_{r_{2}}),
  \label{eq:mutation}
\end{equation}

where $\bm{x}_{r_{1}}$ and $\bm{x}_{r_{2}}$ are randomly selected distinct
individuals. Two hyper-parameters remain fixed throughout the search: the
\CR{} is held constant and determines \Ffactor{} and \GMR{} through
hand-designed equations.
% 待补：说明这两个固定设定为何在本文场景下不够好（可行解比例低、探索-开发失衡）
\todo{explain why these two fixed settings are inadequate here: too few
      feasible solutions and a poor exploration--exploitation balance.}

% 表格占位：基线参数可在此以表格形式给出（可选）
% \input{tables/dmde_parameters}
